\section{A support criterion for arbitrary classical channels}\label{sec:support}
The symmetric family is not required for the boundary value one. We now let $K=(K(y\mid i))$ be any finite column-stochastic matrix with $n\geq2$ inputs. Its quantum extension first dephases $A$ and retains the same finite quantum reference as before. Put
\[
 d(K)=\sum_y\min_i K(y\mid i).
\]
This is the classical Doeblin common-component coefficient. The decomposition giving the upper bound below is prior~\cite{GHNW,HGNW}; the converse is the issue addressed here.

\begin{theorem}[Strict complete contraction and common support]\label{thm:support}
For every such $K$, with the faithful equal-reference Umegaki domain of~\eqref{eq:complete},
\[
 \eta_D^c(K)<1\quad\Longleftrightarrow\quad d(K)>0.
\]
Equivalently, $\eta_D^c(K)=1$ exactly when each nonzero output row has a zero entry. The same criterion holds for $\eta_{D,2\mathrm m}^c(K)$, whose denominator remains quantum. When $d(K)=0$, dimension $n$ of the reference and separable CQ hypotheses with uniform classical marginal suffice. The hypotheses can be chosen independently of $K$: for fixed $n$, their input, output and binary-measured relative entropies tend to $\log2$, uniformly over all finite $K$ with $n$ inputs and $d(K)=0$. The output measurement is allowed to depend on $K$ and the preparation parameter.
\end{theorem}

A common row of positive entries is a stronger support condition than pairwise column overlap. The latter characterizes strict \emph{ordinary} relative-entropy contraction through scrambling~\cite[Theorems~4.1--4.2]{Cohen}. Thus the theorem distinguishes two support tests, rather than claiming a new definition of complete contraction or of Doeblin minorization. For $n=2$ the tests coincide. No claim of strict quantum-reference enhancement is made for channels whose ordinary coefficient is already one.

\subsection{The known upper and the proper-support inverse}
If $0<d<1$, set $a_y=\min_iK(y\mid i)$ and write
\[
 K=dQ+(1-d)L,\qquad Q(y\mid i)=a_y/d,
\]
where $L$ is stochastic and $Q$ is replacement. Equal reference marginals give $Q\rho=Q\sigma$. Joint convexity and data processing imply
\begin{equation}\label{eq:doeblin-upper}
 D(K\rho\Vert K\sigma)\leq(1-d)D(L\rho\Vert L\sigma)
 \leq(1-d)D(\rho\Vert\sigma).
\end{equation}
For $d=1$, $K$ itself is replacement and the coefficient is zero. Consequently $d>0$ implies strict contraction for every finite reference, including the binary-measured coefficient. General quantum inputs reduce to CQ by dephasing, as in Section~2.

For the converse use exactly the hypotheses~\eqref{eq:blocks} and scale~\eqref{eq:scaling}; these also define valid states for $n=2$, with no inactive input modes. Their construction and their input divergence do not depend on $K$. The input calculation in Theorem~\ref{thm:fixed} proves $D_i\to\log2$ also for $n=2$; the exclusion-specific output calculation is not needed.

Discard an identically zero output row only as a common null support. If $d=0$, every remaining row $w_i=K(y\mid i)$ has a proper support $S$, with $1\leq m=|S|\leq n-1$. Write
\[
 t=\sum_iw_i>0,\quad V=\operatorname{span}\{\psi_i:i\in S\},\quad
 T=\sum_iw_i\Pi_i,\quad z=\sum_iw_i\psi_i.
\]
All inverses involving $T$ below are restricted to $V$.
\begin{lemma}\label{lem:support-inverse}
The space $V$ has dimension $m$, $T|_V>0$, and
\[
 z^*T_V^{-1}z=t,\qquad
 \|T_V^{-1}z\|^2=\frac{m(n-1)}{n-m}\leq(n-1)^2.
\]
\end{lemma}
\begin{proof}
Let $B$ have columns $\psi_i$, $i\in S$, and $W=\diag(w_i:i\in S)>0$. The Gram matrix is
\[
 G=B^*B=\frac{nI_m-\mathbf1_m\mathbf1_m^*}{n-1}>0.
\]
Its eigenvalue on $\mathbf1_m$ is $(n-m)/(n-1)$, and its other eigenvalues are $n/(n-1)$. Thus $B$ has full column rank. Since $T=BWB^*$ and $z=BW\mathbf1_m$, the vector $v=BG^{-1}\mathbf1_m$ satisfies $Tv=z$. It follows that $z^*v=t$ and $\|v\|^2=\mathbf1_m^*G^{-1}\mathbf1_m=m(n-1)/(n-m)$. In particular this norm bound has no reciprocal minimum weight.
\end{proof}

\subsection{A controlled eigenvalue, with bounds independent of the row}
On $V\oplus\mathbb Cf$, the minus output block is
\[
 M_- = \begin{pmatrix}A&-g_y\\-g_y^*&b\end{pmatrix},\quad
 A=(\alpha-\eps)T+\eps t I_V,\quad g_y=cz,\quad b=\beta t.
\]
On $P\ominus V$, both output blocks equal $\eps t I$. These common modes remain in the physical state. Since $T\leq tI_V$, we have $A\geq\alpha T$ and $A>\eps t I_V$. Also $\alpha-\eps>0$ for~\eqref{eq:scaling}.

For $0\leq x\leq\eps t$, define the Schur function
\[
 H(x)=b-x-g_y^*(A-xI_V)^{-1}g_y.
\]
Lemma~\ref{lem:support-inverse} yields
\begin{align*}
 H(0)&\geq\beta t\delta>0,\\
 H(\eps t)&=-\frac{t\beta\eps(1-\delta)}{\alpha-\eps}<0,\\
 H'(x)&=-1-\|(A-xI_V)^{-1}g_y\|^2<0.
\end{align*}
There is a unique root $\ell_y\in(0,\eps t)$. Interlacing places the other active eigenvalues above $\eps t$; inactive eigenvalues equal $\eps t$. Therefore $\ell_y$ is the unique global minimum eigenvalue of this nonzero output block for the stated scale. This conclusion is proved here, rather than inferred from a numerical selector.

The operators $A-xI_V$ and $T$ commute, and
\[
 A-xI_V=(\alpha-\eps)T+(\eps t-x)I_V.
\]
Consequently
\[
 \|(A-xI_V)^{-1}g_y\|^2\leq C_{nk},\qquad
 C_{nk}=\frac{c^2(n-1)^2}{(\alpha-\eps)^2}.
\]
Integrating $-H'$ from zero to $\ell_y$ gives
\begin{equation}\label{eq:ell-bound}
 \frac{\beta t\delta}{1+C_{nk}}\leq\ell_y<\beta t\delta.
\end{equation}
For a normalized minimum eigenvector $u_y=(x_y,s_y)$ of $M_-$, with $s_y>0$, the eigenvalue equation gives
\[
 x_y=(A-\ell_y I_V)^{-1}g_y s_y,\qquad
 \frac1{1+C_{nk}}\leq s_y^2\leq1.
\]
These estimates are independent of $S$, the positive row weights and the number of output rows. They do not assert a well-conditioned eigenvector computation.

\subsection{One binary readout and the uniform limit}
Choose $E_y=u_yu_y^*$ on $V\oplus\mathbb Cf$, extended by zero on the inactive subspace, and measure $E=\diag_yE_y$, $I-E$. This is a two-outcome measurement implemented by classical feedforward to a reference measurement. It uses the known hypotheses and $K$, not the unknown sign. For the orientation plus versus minus,
\[
 q_y=\Tr(E_yM_-)=\ell_y,\qquad
 p_y-q_y=4s_y^2g_y^*(A-\ell_y I_V)^{-1}g_y
          =4s_y^2(\beta t-\ell_y).
\]
Summing uses only $\sum_yt_y=n$. Thus, with $p=\sum_yp_y$, $q=\sum_yq_y$,
\begin{align}\label{eq:uniform-pq}
 q_L=\frac{n\beta\delta}{1+C_{nk}}&\leq q\leq q_U=n\beta\delta,\\
 p_L=q_L+\frac{4n\beta(1-\delta)}{1+C_{nk}}
 &\leq p\leq p_U=q_U+4n\beta.\notag
\end{align}
The complement probabilities are $1-p$, $1-q$ on the full normalized states. There is no active-subspace normalization or postselection.

At fixed $n$, $C_{nk}=O_n(1/k)$. These scalar bounds imply, uniformly in the channel,
\[
 p=k^{-1}(1+O_n(k^{-1})),\qquad
 q=\frac{2^{-k}}{4k}(1+O_n(k^{-1})).
\]
Therefore $\log(p/q)=k\log2+\log4+o(1)$ and $p\log(p/q)\to\log2$, while $(1-p)\log[(1-p)/(1-q)]\to0$. This proves $D_b\to\log2$. The inequalities $D_b\leq D_o\leq D_i$, with $D_i\to\log2$, prove the output limit and both ratios. Data processing supplies the upper bound one, establishing Theorem~\ref{thm:support}.

More explicitly, when $p_L>q_U$, put
\[
 L_{nk}=\frac{D(\mathrm{Ber}(p_L)\Vert\mathrm{Ber}(q_U))}{D_i(n,k)},
\]
and put $L_{nk}=0$ otherwise. Monotonicity of Bernoulli KL in $p$ and $q$ for $p>q$ gives simultaneously
\[
 L_{nk}\leq\frac{D_b(K,k)}{D_i(n,k)}\leq
 \frac{D_o(K,k)}{D_i(n,k)}\leq1,\qquad L_{nk}\longrightarrow1.
\]
Hence the same input preparation works for a sequence $K_k$ with changing supports, weights and finite output cardinalities, provided $n$ is fixed and $d(K_k)=0$. The effects $E_y(K_k,k)$ change with the channel. This is neither a channel-independent effect nor uniformity of measurement precision, computational cost, a spectral floor or growing $n$.

\subsection{A certified asymmetric separation and falsifiers}
For example, consider the fixed channel
\[
 K=\begin{pmatrix}
 0&1/3&2/5\\3/4&0&3/5\\1/4&2/3&0
 \end{pmatrix}.
\]
Its Doeblin coefficient is zero, while its classical Dobrushin coefficient is $3/4$; hence its ordinary Umegaki coefficient is at most $3/4$. Directed integer bounds from~\eqref{eq:uniform-pq} certify $D_b/D_i>0.8860189$ at $k=64$ and $D_b/D_i>0.9328480$ at $k=128$. These certificates establish finite separations, not the general theorem.

The rule ``choose the eigenvector of maximal overlap with $f$'' is not a universal finite-$k$ implementation of the minimum eigenprojector. In exclusion with $n=k=16$ it chooses a different branch: the measured ratios are about $0.006905$ for that rule and $0.145103$ for the correct minimum. Both original numerical controls and this falsifier are retained. The proof and readout definition use the minimum established by~\eqref{eq:ell-bound}.

The scale condition is also substantive. In a four-input cycle with two inputs in each output support, replacing $\eps=\beta\delta$ by $\eps=\beta\delta^{1/2}$ changes the measured limiting retention to $1/2$. Replacing it by $\eps=\beta\delta^2$ can instead make a common inactive mode the global minimum; measuring that mode gives zero divergence. These tests do not refute the sealed scale and do not justify an arbitrary small-regularizer substitution.

Finally, the theorem does not give $\eta_D^c=1-d$ away from $d=0$. The binary channel with columns $(0.9,0.1)$ and $(0.2,0.8)$ has complete coefficient $1/2$, whereas $1-d=7/10$. The known upper can be strict. Likewise, appending a marked replacement outcome of mass $0.1$ to a zero-Doeblin channel invalidates measuring a near-$f$ event on every output: that common outcome masks the rare probability. Setting its effect to zero retains its entire mass in the complement and gives a measured limit $0.9\log2$ for this marked example. The failed first protocol remains in the supplement. No formula for a general positive-$d$ optimum, robust communication advantage, capacity or optimal copy complexity is inferred.
